\documentclass{article}
\usepackage{amsmath, amssymb, amsthm}
\usepackage{hyperref}
\usepackage{geometry}
\geometry{margin=1in}

\title{Erd\H{o}s Problem \#790}
\date{Last updated: 2025-08-31}
\author{Source: erdosproblems.com}

\begin{document}
\maketitle

\noindent\textbf{Status:} OPEN \quad \textbf{No prize}

\noindent\textbf{Tags:} additive combinatorics

\medskip

\noindent\textbf{Problem Statement:}

Let $l(n)$ be maximal such that if $A\subset\mathbb{Z}$ with $\lvert A\rvert=n$ then there exists a sum-free $B\subseteq A$ with $\lvert B\rvert \geq l(n)$ - that is, $B$ is such that there are no solutions to\[a_1=a_2+\cdots+a_r\]with $a_i\in B$ all distinct. Estimate $l(n)$. In particular, is it true that $l(n)n^{-1/2}\to \infty$? Is it true that $l(n)&lt; n^{1-c}$ for some $c&gt;0$?




Erd\H{o}s observed that $l(n)\geq (n/2)^{1/2}$, which Choi improved to $l(n)&gt;(1+c)n^{1/2}$ for some $c&gt;0$. Erd\H{o}s \cite{Er73} thought he could prove $l(n)=o(n)$ but had 'difficulties in reconstructing [his] proof'. (In \cite{Er65} he wrote 'by complicated arguments we can show $l(n)=o(n)$'.) Choi, Koml\'{o}s, and Szemer\'{e}di \cite{CKS75} proved\[\left(\frac{\log n}{\log\log n}n\right)^{1/2}\ll l(n) \ll \frac{n}{\log n}.\]They further conjecture that $l(n)\geq n^{1-o(1)}$. See also [876] .




References

[CKS75] Choi, S. L. G. and Koml\'os, J. and Szemer\'{e}di, E., On sum-free subsequences . Trans. Amer. Math. Soc. (1975), 307--313.

[Er65] Erd\H{o}s, P., Extremal problems in number theory . Proc. Sympos. Pure Math., Vol. VIII (1965), 181-189.

[Er73] Erd\H{o}s, P., Problems and results on combinatorial number theory . A survey of combinatorial theory (Proc. Internat. Sympos., Colorado State Univ., Fort Collins, Colo., 1971) (1973), 117-138.

\noindent\textbf{Related OEIS sequences:} possible


\bigskip
\noindent\small{Source: \url{https://www.erdosproblems.com/790}}
\end{document}
